(1) 每相满足 \((\partial S/\partial T)_p=C_p/T\)。若 \(\Delta V\) 与温度无关，Maxwell关系给出 \(\partial(\Delta S)/\partial p=-(\partial\Delta V/\partial T)_p=0\)，因此 \(\Delta S(T)=L_0/T_0+\Delta C_p\ln(T/T_0)\)，该表达式沿共存线同样成立。
(2) Clapeyron式 \(dp/dT=\Delta S/\Delta V\)。从 \(T_0\) 积分得 \(p(T)=p_0+\frac{L_0}{T_0\Delta V}(T-T_0)+\frac{\Delta C_p}{\Delta V}[T\ln(T/T_0)-T+T_0]\)，直接求导可恢复前述斜率。
(3) 相变潜热沿共存线满足 \(L(T)=T\Delta S(T)=T[L_0/T_0+\Delta C_p\ln(T/T_0)]\)；共存曲率 \(d^2p/dT^2=\Delta C_p/(\Delta VT)\)。注意潜热不等于简单的 \(L_0+\Delta C_p(T-T_0)\)，因为沿共存线压强也改变。
(4) 当 \(\Delta C_p\to0\)，\(p-p_0=L_0(T-T_0)/(T_0\Delta V)\)，曲线线性。若 \(\Delta V\) 实际随压温变化显著，或接近临界点，恒体积差假设失效，需恢复一般Clapeyron积分而不可外推此解析式。
